\begin{lemma}{Integrierbarkeitskriterien}
             {Integrierbarkeitskriterien}
Seien $f, g, h \in \mcm(\Omega, \mfa, \R)$.\\
Dann gilt:
\begin{enumerate}[(1)]
	\item \label{Integrierbarkeitskriterien1}
 Sind $f, g \ge 0$, $f \le g$ und $g$ integrierbar, so auch $f$ mit
$0 \le \E(f) \le \E(g)$.
\item \label{Integrierbarkeitskriterien2}
 Sind $h \le f \le g$ und $g, h$ integrierbar, so auch $f$ und
$\E(h) \le \E(f) \le \E(g)$.
\end{enumerate}
\end{lemma}
\begin{proof}
Ist die Integrierbarkeit von $f$ in den beiden F�llen gekl�rt, so folgt die
jeweils zweite Aussage einfach �ber die Monotonie des Erwartungswerts.\\
(1) Seien $f, g \ge 0$, $f\le g$ und $g$ integrierbar.\\
Sei $\vi \in \mct(\Omega, \mfa, \R)$ mit $0 \le \vi \le f$.\\
Dann ist auch $0 \le \vi \le f \le g$ und daher
\[\menge{\E(\psi)}{\psi \in \mct(\Omega,\mfa,\R), 0 \le \psi\le f}
\subseteq
 \menge{\E(\psi)}{\psi \in \mct(\Omega,\mfa,\R), 0 \le \psi\le g}\text.\]
Die gr��ere Menge ist allerdings beschr�nkt, weil $g$ integrierbar ist.\\
Also ist $\menge{\E(\psi)}{\psi \in \mct_1(\Omega,\mfa,\R), 0 \le \psi\le f}$
nichtleer und nach oben beschr�nkt, besitzt also ein Supremum in $\R$.\\
Daher
\begin{align*}
\E(f)&=\sup\menge{\E(\psi)}{\psi \in \mct(\Omega,\mfa,\R), 0 \le \psi\le f}\\
   &\le\sup\menge{\E(\psi)}{\psi \in \mct(\Omega,\mfa,\R), 0 \le \psi\le g}
   = \E(g) < \infty\text.
\end{align*}
Somit ist $f$ integrierbar.\\
(2) Seien $h \le f \le g$ und $g, h$ integrierbar.\\
Es gilt also $0 \le f-h \le g -h$.\\
Da $g, h$ integrierbar sind, ist es auch $g-h$ und damit nach (1) auch $f-h$.\\
Analog ist also $(f-h)+ h = f$ integrierbar.
\end{proof}